\documentclass[12pt]{article}
\usepackage[T1]{fontenc}
\usepackage[utf8]{inputenc}
\usepackage{lmodern}
\usepackage{textcomp}
\usepackage{enumitem}
\usepackage{geometry}
\geometry{margin=1in}
\begin{document}
\begin{center}
Answer Key - Machine Learning - Undergraduate Level Assessment 8 \
\small (Answers are ordered exactly as in the question paper.)
\end{center}

\section*{Multiple Choice}
\begin{enumerate}[label=\arabic*.]
\item \textbf{Answer:} B
\\ \textit{Options:} A) Larger if the error rate is larger.; B) Larger if the error rate is smaller.; C) Smaller if the error rate is smaller.; D) It does not matter.
\\ \textit{Note:} Correct option: B - Larger if the error rate is smaller.
\item \textbf{Answer:} B
\\ \textit{Options:} A) Increase the amount of training data.; B) Improve the optimisation algorithm being used for error minimisation.; C) Decrease the model complexity.; D) Reduce the noise in the training data.
\\ \textit{Note:} Correct option: B - Improve the optimisation algorithm being used for error minimisation.
\item \textbf{Answer:} A
\\ \textit{Options:} A) random crop and horizontal flip; B) random crop and vertical flip; C) posterization; D) dithering
\\ \textit{Note:} Correct option: A - random crop and horizontal flip
\item \textbf{Answer:} A
\\ \textit{Options:} A) The polynomial degree; B) Whether we learn the weights by matrix inversion or gradient descent; C) The assumed variance of the Gaussian noise; D) The use of a constant-term unit input
\\ \textit{Note:} Correct option: A - The polynomial degree
\item \textbf{Answer:} D
\\ \textit{Options:} A) It can not be applied to non-differentiable functions.; B) It can not be applied to non-continuous functions.; C) It is hard to implement.; D) It runs reasonably slow for multiple linear regression.
\\ \textit{Note:} Correct option: D - It runs reasonably slow for multiple linear regression.
\end{enumerate}

\section*{True / False}
\begin{enumerate}[label=\arabic*.]
\setcounter{enumi}{5}
\item \textbf{Answer:} False
\\ \textit{Options:} A) True; B) False
\\ \textit{Note:} LLM-generated false statement. Correct answer: 'P(X, Y, Z) = P(Y) * P(X|Y) * P(Z|Y)'.
\item \textbf{Answer:} False
\\ \textit{Options:} A) True; B) False
\\ \textit{Note:} LLM-generated false statement. Correct answer: 'True, True'.
\item \textbf{Answer:} True
\\ \textit{Options:} A) True; B) False
\\ \textit{Note:} LLM-generated true statement based on MCQ.
\item \textbf{Answer:} True
\\ \textit{Options:} A) True; B) False
\\ \textit{Note:} LLM-generated true statement based on MCQ.
\item \textbf{Answer:} True
\\ \textit{Options:} A) True; B) False
\\ \textit{Note:} LLM-generated true statement based on MCQ.
\end{enumerate}

\section*{Long Form Response}
\begin{enumerate}[label=\arabic*.]
\setcounter{enumi}{10}
\item \textbf{Answer:} def first\_repeated\_word(str 1): | return word; | return 'None'
\\ \textit{Note:} Students should provide a detailed explanation referencing algorithm steps.
\item \textbf{Answer:} def count\_vowels(test\_str): | return (res)
\\ \textit{Note:} Students should provide a detailed explanation referencing algorithm steps.
\end{enumerate}

\vfill
\noindent\textit{End of Answer Key}
\end{document}